% Source: Weinberg GR, physical PDF pp. 205--208
%         (printed pp. 182--185).
% Coverage: Section 8.3, Other Metrics; equations (8.3.1)--(8.3.9).
% Figures/tables/footnotes: reference marker 1.
% Status: source-reviewed and compile-clean.
% Uncertainties: none.

\subsection{Other Metrics}\label{sec:8.3}

The general kinematic framework provided by the Principle of Equivalence
rests on a much firmer foundation than do Einstein's field equations.
Indeed, in Chapters~3 through~5 we were led almost inevitably from the
equality of gravitational and inertial mass to the full formalism of tensor
analysis and general covariance, whereas in contrast the derivation of
Einstein's equations in Chapter~7 contained a strong element of guesswork,
and in any case there might exist a long-range scalar field, like that of
Brans and Dicke, that would alter the field equations. It is therefore very
useful to test general relativity by assuming that the usual rules for the
motion of particles and photons in a given metric field \(g_{\mu\nu}\) still
apply, but that the metric may be different from that calculated from the
Einstein equations.

In any case we would expect the metric produced by a static spherically
symmetric body like the sun to be expressible in the ``standard,''
``isotropic,'' and ``harmonic'' forms given in Section~\ref{sec:8.1}, and
we would further expect that the metric coefficients [e.g., \(A(r)\) and
\(B(r)\)] could be expanded as power series in the small parameter \(MG/r\).
Such an expansion was given by Eddington and
Robertson\hyperref[ch8-ref-1]{\textsuperscript{1}} for the metric in its
isotropic form:
\begin{equation}
  \begin{split}
  \dd s^2={}&
  -\left(
    1-2\alpha\frac{MG}{\rho}
    +2\beta\frac{M^2G^2}{\rho^2}+\cdots
  \right)\dd t^2\\
  &+
  \left(1+2\gamma\frac{MG}{\rho}+\cdots\right)
  \left(\dd\rho^2+\rho^2\dd\theta^2
  +\rho^2\sin^2\theta\,\dd\varphi^2\right),
  \end{split}
  \tag{8.3.1}\label{eq:8.3.1}
\end{equation}
where \(\alpha,\beta\), and \(\gamma\) are unknown dimensionless
parameters. (The reason for carrying this expansion to order
\(M^2G^2/\rho^2\) in \(g_{00}\) and only to order \(MG/\rho\) in
\(g_{ij}\) is that in applications to celestial mechanics \(g_{ij}\) will
always get multiplied with an extra factor \(v^2\sim MG/\rho\).)
Comparing with the isotropic form \eqref{eq:8.2.14} of the Schwarzschild
solution, we see that the predictions of the Einstein field equations can
be neatly summarized as
\begin{equation}
  \alpha=\beta=\gamma=1.
  \tag{8.3.2}\label{eq:8.3.2}
\end{equation}
In contrast, the Brans--Dicke theory discussed in Section~\ref{sec:7.3}
gives a metric (see Section~\ref{sec:9.9}) that can be expressed as in
Eq.~\eqref{eq:8.3.1}, with
\begin{equation}
  \alpha=\beta=1,\qquad
  \gamma=\frac{\omega+1}{\omega+2},
  \tag{8.3.3}\label{eq:8.3.3}
\end{equation}
where \(\omega\) is the unknown dimensionless parameter of this theory. In
order to decide whether Einstein, or Brans and Dicke, or someone else, has
the right field equations, what must be done is to measure
\(\alpha,\beta\), and \(\gamma\).

We shall generally be doing our calculations with the metric in its
``standard'' form, so it will prove convenient to convert the Robertson
expansion \eqref{eq:8.3.1} to this form by defining
\begin{equation}
  r\equiv\rho\left(1+\gamma\frac{MG}{\rho}+\cdots\right),
  \tag{8.3.4}\label{eq:8.3.4}
\end{equation}
or
\[
  \rho=r\left(1-\gamma\frac{MG}{r}+\cdots\right).
\]
A simple calculation gives
\begin{equation}
  \begin{split}
  \dd s^2={}&
  -\left[
    1-2\alpha\frac{MG}{r}
    +2(\beta-\alpha\gamma)\frac{M^2G^2}{r^2}
    +\cdots
  \right]\dd t^2\\
  &+
  \left(1+2\gamma\frac{MG}{r}+\cdots\right)\dd r^2
  +r^2\dd\theta^2+r^2\sin^2\theta\,\dd\varphi^2.
  \end{split}
  \tag{8.3.5}\label{eq:8.3.5}
\end{equation}

Also, we can construct harmonic coordinates \(\mathbf{X},t\), by using
\[
  X_1=R\sin\theta\cos\varphi,\qquad
  X_2=R\sin\theta\sin\varphi,\qquad
  X_3=R\cos\theta,
\]
with \(R\) satisfying the differential equation \eqref{eq:8.1.15}:
\[
  0
  =
  \frac{\dd}{\dd r}
  \left[
    r^2\left(1-(\alpha+\gamma)\frac{MG}{r}+\cdots\right)
    \frac{\dd R}{\dd r}
  \right]
  -2\left(1-(\alpha-\gamma)\frac{MG}{r}+\cdots\right)R.
\]
The solution is
\begin{equation}
  R
  =
  \left(1+\frac{(\alpha-3\gamma)MG}{2r}+\cdots\right)r,
  \tag{8.3.6}\label{eq:8.3.6}
\end{equation}
and Eq.~\eqref{eq:8.1.16} gives the metric (with
\(R^2\equiv\mathbf{X}^2\))
\begin{equation}
  \begin{split}
  \dd s^2={}&
  -\left[
    1-2\alpha\frac{MG}{R}
    +(\alpha\gamma-\alpha^2+2\beta)\frac{M^2G^2}{R^2}
    +\cdots
  \right]\dd t^2\\
  &+
  \left[1+\frac{(3\gamma-\alpha)MG}{R}+\cdots\right]
  \dd\mathbf{X}^2\\
  &+
  \frac{(\alpha-\gamma)MG/R+\cdots}{R^2}
  (\mathbf{X}\cdot \dd\mathbf{X})^2.
  \end{split}
  \tag{8.3.7}\label{eq:8.3.7}
\end{equation}
Comparing Eqs.~\eqref{eq:8.3.5} and \eqref{eq:8.3.7} with the
corresponding exact solutions \eqref{eq:8.2.12} and \eqref{eq:8.2.15}
shows again that Einstein's theory gives \(\alpha=\beta=\gamma=1\).

The prediction that \(\alpha=1\) really just follows from the empirical
definition of the mass \(M\). Note that Eq.~\eqref{eq:8.3.1} would give a
slowly moving particle far from the origin a centripetal acceleration equal
to
\[
  -g=-\Gamma^r{}_{tt}
  =\frac12\partial_r g_{tt}
  =-\frac{\alpha MG}{r^2}
  \qquad
  (MG/r\ll1\ \hbox{and}\ v^2\ll1),
\]
whereas in fact the masses of the sun and planets are measured by setting
\(g=MG/r^2\); hence we must absorb \(\alpha\) into \(M\), or, in other
words, we must choose \(\alpha=1\). Only if it were possible to determine
\(M\) by some independent nongravitational measurement would it make sense
to ask whether in fact \(\alpha\) is exactly unity.

With \(\alpha=1\), the metric functions given by Eq.~\eqref{eq:8.3.5}
are
\begin{equation}
  B(r)
  =
  1-\frac{2MG}{r}
  +2(\beta-\gamma)\frac{M^2G^2}{r^2}+\cdots,
  \tag{8.3.8}\label{eq:8.3.8}
\end{equation}
\begin{equation}
  A(r)=1+2\gamma\frac{MG}{r}+\cdots.
  \tag{8.3.9}\label{eq:8.3.9}
\end{equation}
As shown in Chapter~3, the gravitational red shift experiment only
measures the term \(-2MG/r\) in \(B(r)\), and hence can only verify the
Principle of Equivalence. We shall see that, of the other tests of general
relativity listed at the beginning of this chapter, (B) and (D) can only
test whether \(\gamma\simeq1\), whereas (C), the precession of perihelia,
verifies that \(2\gamma-\beta\simeq1\). (To the extent that we ignore the
rotation of the earth, (E) also only tests whether \(\gamma\simeq1\).)
